\section{Open systems: escape, amplitudes, and resonances}
\label{sec:open-system}

\subsection{First-entry residue events and escape rates}

We now pass from general boundary transfer to the central open-system
specialization of the paper.  For first-entry residue events, the
boundary matrix is precisely the killed transfer matrix of a symbolic
hole.  This turns the finite-observation problem into a question about
escape, quasistationary amplitudes, and cyclotomic lifts of open
spectra.  The study of escape rates through holes in dynamical systems
originates with
\citet{PianigianiYorke1979}; the transfer-operator approach for
topological Markov chains was developed by
\citet{ColletMartinezSchmitt1997}.  For spectral characterizations of
escape rates and conditionally invariant measures, see
\citet{DemersYoung2006,DemersWrightYoung2012} and
\citet{KellerLiverani2009}; for the Gibbs-measure setting, see
\citet{FergusonPollicott2012}.  Since any finite union of cylinders
becomes a one-step Markov hole after a sufficiently high block
presentation, there is no loss in working directly on a one-step shift
over a finite alphabet \(A\); compare \citet{LindMarcus1995,Bowen2008}.

\begin{theorem}[First-entry residue events and escape rates]
\label{thm:first-entry-escape-rate}
Let \(Y\subseteq A^{\NNzero}\) be a topologically mixing one-step
subshift of finite type, and let \(\nu_\varphi\) be the fully supported
equilibrium Markov measure of a one-step potential \(\varphi\).  Write
\(\pi\) for the stationary vector and \(K=(K_{ij})_{i,j\in A}\) for the
transition matrix.  Let \(M=(M_{ij})_{i,j\in A}\) be the adjacency
matrix of \(Y\).  Let \(A_{\mathrm{hole}}\subsetneq A\) be a non-empty proper
subset, and set
\[
  H:=\bigcup_{h\in A_{\mathrm{hole}}}[h],
  \qquad
  \mathcal S:=A\setminus A_{\mathrm{hole}}.
\]
Let
\[
  Y_H:=\set{y\in Y:\sigma^n y\notin H\text{ for every }n\geq 0}
\]
be the survivor subshift, and assume that \(Y_H\) is topologically
mixing.  Let
\[
  \mathcal S_\infty
  := \set{i\in\mathcal S:\exists\,y\in Y_H,\ y_0=i}
\]
be the survivor alphabet, i.e.\ the set of states that actually appear
in \(Y_H\).  Since \(Y_H\) is a one-step subshift of finite type, its
adjacency matrix restricted to \(\mathcal S_\infty\) is irreducible;
mixing further implies primitivity.  Set
\[
  \mathcal T:=\mathcal S\setminus\mathcal S_\infty.
\]
States in \(\mathcal T\) are transient: from any such state, every
admissible path eventually enters \(A_{\mathrm{hole}}\).
Define the first entry time
\[
  \tau_H(y):=\inf\set{n\geq 0:\sigma^n y\in H}.
\]
Fix \(q\geq 2\) and a non-empty proper subset
\(R\subsetneq \ZZ/q\ZZ\), and define
\[
  \mathsf{P}_R
  := \set{y\in Y:
      \tau_H(y)<\infty,\
      \tau_H(y)\bmod q\in R }.
\]

For \(i\in\mathcal S\), let \(\PP_i\) be the Markov law on \(Y\)
conditioned to start at state \(i\), and let
\[
  \tau_H^+
  := \inf\set{n\geq 1:y_n\in A_{\mathrm{hole}}}.
\]
For \(r\in\ZZ/q\ZZ\), define
\[
  p_{i,r}:=\PP_i\bigl(r+\tau_H^+\in R \bmod q\bigr),
  \qquad
  \beta_{i,r}:=\min\set{p_{i,r},1-p_{i,r}}.
\]
Assume the following residue non-degeneracy condition:
for every \(i\in\mathcal S\) and \(r\in\ZZ/q\ZZ\), there exist
admissible future paths starting at \(i\) whose first entry into
\(A_{\mathrm{hole}}\) occurs at residues both in \(R-r\) and in
\((\ZZ/q\ZZ\setminus R)-r\).  Then there exists \(\theta>0\) such that
\[
  0<\theta\leq \beta_{i,r}\leq \frac12
  \qquad
  \text{for every } i\in\mathcal S,\ r\in\ZZ/q\ZZ.
\]

Define the killed matrix \(B_H\) on \(\mathcal S\) by
\[
  B_H(i,j):=K_{ij}\mathbf{1}_{\{j\in\mathcal S\}},
  \qquad i,j\in\mathcal S,
\]
let \(u\in\RR^{\mathcal S}\) be given by
\[
  u_i:=\pi_i,
\]
and for \(r\in\ZZ/q\ZZ\) let \(b_r\in\RR^{\mathcal S}\) be the vector
\[
  b_r(i):=\beta_{i,r}.
\]
Then for every \(m\geq 1\),
\[
  \varepsilon_m(\mathsf{P}_R;\nu_\varphi)
  = u^{\mathsf T}B_H^{m-1}b_{m-1\bmod q}.
\]
Moreover,
\[
  \theta\,\nu_\varphi(\tau_H\geq m)
  \leq \varepsilon_m(\mathsf{P}_R;\nu_\varphi)
  \leq \frac12\,\nu_\varphi(\tau_H\geq m),
\]
where
\[
  \nu_\varphi(\tau_H\geq m)
  = u^{\mathsf T}B_H^{m-1}\mathbf{1}.
\]
If \(\rho_H:=\rho(B_H)\), then
\[
  \lim_{m\to\infty}
  -\frac1m\log\varepsilon_m(\mathsf{P}_R;\nu_\varphi)
  = -\log\rho_H.
\]
Finally, if
\[
  L_{\varphi,H}(i,j):=M_{ij}e^{\varphi(i,j)}\mathbf{1}_{\{j\in\mathcal S_\infty\}},
  \qquad i,j\in\mathcal S_\infty,
\]
then
\[
  -\log\rho_H
  = P_Y(\varphi)-P_{Y_H}(\varphi).
\]
Thus the finite-observation exponent of \(\mathsf{P}_R\) coincides with
the escape rate of the open system \((Y,\sigma,H)\).
\end{theorem}

\begin{proof}
Let \(a=a_0\cdots a_{m-1}\in A^m\).  If \(a_t\in A_{\mathrm{hole}}\) for some
\(0\leq t\leq m-1\), then \(\tau_H\leq m-1\) on the cylinder \([a]\), so
the truth value of \(\mathsf{P}_R\) is already decided on \([a]\), and
\([a]\) contributes \(0\) to
\(\varepsilon_m(\mathsf{P}_R;\nu_\varphi)\).

Assume therefore that \(a_t\in\mathcal S\) for every
\(0\leq t\leq m-1\).  By the Markov property, conditioned on \([a]\) the
future depends on the prefix only through the terminal state
\(a_{m-1}\).  Since the next possible hit time occurs after at least one
additional step, the conditional probability of \(\mathsf{P}_R\) on
\([a]\) is exactly \(p_{a_{m-1},\,m-1\bmod q}\).  Hence the Bayes factor
on \([a]\) is \(\beta_{a_{m-1},\,m-1\bmod q}\).  The residue
non-degeneracy hypothesis guarantees \(0<\beta_{i,r}\leq \tfrac12\), and
since the state space \(\mathcal S\times \ZZ/q\ZZ\) is finite, the lower
bound can be chosen uniformly.

Thus the boundary cylinders for \(\mathsf{P}_R\) are exactly the
length-\(m\) survivor words, and
Proposition~\ref{prop:scan-error-cylinder} gives
\[
  \varepsilon_m(\mathsf{P}_R;\nu_\varphi)
  = \sum_{\substack{a\in A^m\\ a_t\in\mathcal S\ \forall t<m}}
      \nu_\varphi([a])\,
      \beta_{a_{m-1},\,m-1\bmod q}.
\]
Since \(\nu_\varphi\) is Markov,
\[
  \nu_\varphi([a])
  = \pi_{a_0}K_{a_0a_1}\cdots K_{a_{m-2}a_{m-1}},
\]
and summing over survivor words with a fixed terminal state gives
\[
  \varepsilon_m(\mathsf{P}_R;\nu_\varphi)
  = u^{\mathsf T}B_H^{m-1}b_{m-1\bmod q}.
\]
The same path sum with terminal weight \(\mathbf{1}\) gives
\[
  \nu_\varphi(\tau_H\geq m)
  = u^{\mathsf T}B_H^{m-1}\mathbf{1}.
\]
The two-sided comparison with survivor mass follows immediately from the
uniform bounds on \(\beta_{i,r}\).

Set \(B_\infty:=B_H|_{\mathcal S_\infty}\).  Since \(Y_H\) is
topologically mixing, \(B_\infty\) is primitive.  Reorder the states in
\(\mathcal S\) so that those in \(\mathcal S_\infty\) come first and
those in \(\mathcal T\) come last.  Then
\[
  B_H=
  \begin{pmatrix}
    B_\infty & C\\
    0 & Q
  \end{pmatrix}.
\]
There is no edge from \(\mathcal T\) into \(\mathcal S_\infty\), which
explains the zero block: if \(i\in\mathcal T\) and \(B_H(i,j)>0\) for
some \(j\in\mathcal S_\infty\), then following an infinite survivor path
from \(j\) would produce an element of \(Y_H\) starting at \(i\), a
contradiction.  Moreover the directed graph on \(\mathcal T\) is
acyclic.  Indeed, a directed cycle inside \(\mathcal T\) could be
repeated indefinitely to produce an infinite admissible path avoiding
\(H\), so its initial state would belong to \(\mathcal S_\infty\), again
a contradiction.  Since \(\mathcal T\) is finite, acyclicity implies
that \(Q\) is nilpotent.  Hence
\[
  \rho_H=\rho(B_H)=\max\set{\rho(B_\infty),\rho(Q)}=\rho(B_\infty).
\]

Let \(u_\infty\) and \(b_{r,\infty}\) denote the restrictions of \(u\)
and \(b_r\) to \(\mathcal S_\infty\).  These vectors are strictly
positive because \(\pi_i>0\) for every \(i\in A\) and
\(\beta_{i,r}\geq \theta\).  Perron--Frobenius therefore gives, for each
\(r\in\ZZ/q\ZZ\),
\[
  u_\infty^{\mathsf T}B_\infty^{m-1}b_{r,\infty}
  = c_r\rho_H^{m-1}+O(\vartheta^m)
\]
with \(c_r>0\) and \(0\leq \vartheta<\rho_H\).  Since all matrices and
vectors are non-negative,
\[
  u^{\mathsf T}B_H^{m-1}b_r
  \geq
  u_\infty^{\mathsf T}B_\infty^{m-1}b_{r,\infty}.
\]
On the other hand, for every \(\epsilon>0\), finite-dimensional spectral
theory yields a constant \(C_\epsilon>0\) such that
\[
  \norm{B_H^{m-1}}\leq C_\epsilon(\rho_H+\epsilon)^{m-1}
  \qquad (m\geq 1).
\]
Because \(q\) is finite, the vectors \(b_r\) take only finitely many
values.  Combining the last two displays, and then taking \(r=m-1\bmod
q\), shows that
\[
  c\,\rho_H^{m-1}
  \leq u^{\mathsf T}B_H^{m-1}b_{m-1\bmod q}
  \leq C_\epsilon'(\rho_H+\epsilon)^{m-1}
\]
for suitable constants \(c>0\) and \(C_\epsilon'>0\).  Hence
\[
  \lim_{m\to\infty}
  -\frac1m\log\varepsilon_m(\mathsf{P}_R;\nu_\varphi)
  = -\log\rho_H.
\]

Let \(h>0\) be a Perron right eigenvector of the ambient transfer matrix
\(L_\varphi(i,j):=M_{ij}e^{\varphi(i,j)}\).  The normalized transition
matrix satisfies
\[
  K_{ij}
  = M_{ij}e^{\varphi(i,j)-P_Y(\varphi)}\frac{h_j}{h_i}.
\]
If \(D\) denotes the diagonal matrix on \(\mathcal S_\infty\) with
\(D_{ii}=h_i\), then the restriction to \(\mathcal S_\infty\) satisfies
\[
  B_H|_{\mathcal S_\infty}=e^{-P_Y(\varphi)}D^{-1}L_{\varphi,H}D.
\]
Consequently
\[
  \rho_H
  = \rho(B_H|_{\mathcal S_\infty})
  = e^{-P_Y(\varphi)}\rho(L_{\varphi,H}).
\]
Since \(Y_H\) is a topologically mixing one-step subshift with alphabet
\(\mathcal S_\infty\), standard thermodynamic formalism identifies
\(\log\rho(L_{\varphi,H})\) with \(P_{Y_H}(\varphi)\).  Therefore
\[
  -\log\rho_H
  = P_Y(\varphi)-P_{Y_H}(\varphi).
\]
The escape-rate interpretation follows from the comparison with
\(\nu_\varphi(\tau_H\geq m)\).
\end{proof}

\begin{corollary}[Survivor pressure recovery]
\label{cor:survivor-pressure-recovery}
Under the hypotheses of
Theorem~\ref{thm:first-entry-escape-rate},
\[
  P_{Y_H}(\varphi)
  = P_Y(\varphi)
    + \lim_{m\to\infty}\frac1m
      \log \varepsilon_m(\mathsf{P}_R;\nu_\varphi).
\]
In particular, for the Parry measure of maximal entropy
\cite{Parry1964},
\[
  h_{\mathrm{top}}(Y_H)
  = h_{\mathrm{top}}(Y)
    + \lim_{m\to\infty}\frac1m
      \log \varepsilon_m(\mathsf{P}_R;\nu_{\max}).
\]
\end{corollary}

\begin{proof}
This is a rearrangement of the pressure identity in
Theorem~\ref{thm:first-entry-escape-rate}.
\end{proof}

\begin{theorem}[Quasistationary ambiguity amplitudes]
\label{thm:quasistationary-ambiguity}
Under the hypotheses of Theorem~\ref{thm:first-entry-escape-rate},
choose an ordering of \(\mathcal S\) for which
\[
  B_H=
  \begin{pmatrix}
    B_\infty & C\\
    0 & Q
  \end{pmatrix},
\]
where \(B_\infty=B_H|_{\mathcal S_\infty}\) and \(Q\) is nilpotent.  Let
\(r_H\) and \(\ell_\infty\) be the Perron right and left eigenvectors of
\(B_\infty\), normalized by
\[
  B_\infty r_H=\rho_H r_H,
  \qquad
  \ell_\infty^{\mathsf T}B_\infty=\rho_H\ell_\infty^{\mathsf T},
  \qquad
  \ell_\infty^{\mathsf T}r_H=1.
\]
Define a row vector on the whole safe alphabet by
\[
  \ell_H^{\mathsf T}
  :=
  \bigl(
    \ell_\infty^{\mathsf T},
    \ell_\infty^{\mathsf T}C(\rho_H I-Q)^{-1}
  \bigr),
  \qquad
  \xi_H:=\frac{\ell_H}{\ell_H^{\mathsf T}\mathbf 1}.
\]
Then \(\ell_H^{\mathsf T}B_H=\rho_H\ell_H^{\mathsf T}\), and for every
\(r\in\ZZ/q\ZZ\),
\[
  \lim_{n\to\infty}
  \frac{
    \varepsilon_{nq+r+1}(\mathsf P_R;\nu_\varphi)
  }{
    \nu_\varphi(\tau_H\ge nq+r+1)
  }
  = \xi_H^{\mathsf T}b_r.
\]
Equivalently,
\[
  \varepsilon_m(\mathsf P_R;\nu_\varphi)
  = \bigl(\xi_H^{\mathsf T}b_{m-1\bmod q}+o(1)\bigr)\,
    \nu_\varphi(\tau_H\ge m).
\]
\end{theorem}

\begin{proof}
Since \(Q\) is nilpotent, there exists \(M\geq 1\) such that \(Q^M=0\).
Let \(u_\infty\) be the restriction of \(u\) to \(\mathcal S_\infty\).
For every \(n\geq 1\),
\[
  B_H^n=
  \begin{pmatrix}
    B_\infty^n &
    \sum_{j=0}^{n-1}B_\infty^{n-1-j}CQ^j\\
    0 & Q^n
  \end{pmatrix},
\]
and therefore for \(n\geq M\) and every vector
\(v=(v_\infty,v_T)\in\RR^{\mathcal S_\infty}\times\RR^{\mathcal T}\),
\[
  u^{\mathsf T}B_H^n v
  =
  u_\infty^{\mathsf T}B_\infty^n v_\infty
  + \sum_{j=0}^{M-1}
      u_\infty^{\mathsf T}B_\infty^{n-1-j}CQ^j v_T.
\]

By Perron--Frobenius,
\[
  B_\infty^n
  = \rho_H^n r_H\ell_\infty^{\mathsf T}+O(\vartheta^n)
\]
for some \(0\leq \vartheta<\rho_H\).  Substituting this expansion into
the previous display gives
\begin{align*}
  u^{\mathsf T}B_H^n v
  &=
    (u_\infty^{\mathsf T}r_H)\rho_H^n
    \left(
      \ell_\infty^{\mathsf T}v_\infty
      + \sum_{j=0}^{M-1}
          \rho_H^{-j-1}\ell_\infty^{\mathsf T}CQ^j v_T
    \right)
    + O(\vartheta^n).
\end{align*}
Because \(Q^M=0\),
\[
  (\rho_H I-Q)^{-1}
  = \sum_{j=0}^{M-1}\rho_H^{-j-1}Q^j.
\]
Hence the coefficient in parentheses is exactly \(\ell_H^{\mathsf T}v\),
and
\[
  u^{\mathsf T}B_H^n v
  =
  (u_\infty^{\mathsf T}r_H)(\ell_H^{\mathsf T}v)\rho_H^n
  + O(\vartheta^n).
\]
Taking \(v=b_r\) and using
Theorem~\ref{thm:first-entry-escape-rate} with \(m=nq+r+1\) gives
\[
  \varepsilon_{nq+r+1}(\mathsf P_R;\nu_\varphi)
  =
  (u_\infty^{\mathsf T}r_H)(\ell_H^{\mathsf T}b_r)\rho_H^{nq+r}
  + O(\vartheta^{nq+r}).
\]
Taking \(v=\mathbf 1\) gives
\[
  \nu_\varphi(\tau_H\ge nq+r+1)
  =
  (u_\infty^{\mathsf T}r_H)(\ell_H^{\mathsf T}\mathbf 1)\rho_H^{nq+r}
  + O(\vartheta^{nq+r}).
\]
Since \(u_\infty>0\), \(r_H>0\), and \(\ell_H^{\mathsf T}\mathbf 1>0\),
the ratio tends to
\[
  \frac{\ell_H^{\mathsf T}b_r}{\ell_H^{\mathsf T}\mathbf 1}
  = \xi_H^{\mathsf T}b_r.
\]
The equivalent statement is immediate on writing \(m-1=nq+r\).
\end{proof}

\begin{proof}[Proof of Theorem~\ref{thm:main-amplitude}]
This is exactly Theorem~\ref{thm:quasistationary-ambiguity}.
\end{proof}

\begin{theorem}[Error resolvent and cyclotomic resonance lift]
\label{thm:error-resolvent-cyclotomic}
Under the hypotheses of
Theorem~\ref{thm:first-entry-escape-rate}, set
\[
  \widehat{\mathcal S}:=\mathcal S\times \ZZ/q\ZZ,
\]
and define the lifted killed matrix \(\widehat B_H\) on
\(\widehat{\mathcal S}\) by
\[
  \widehat B_H\bigl((i,r),(j,s)\bigr)
  := B_H(i,j)\mathbf{1}_{\{s=r+1\bmod q\}}.
\]
Let \(\widehat u,\widehat b\in\RR^{\widehat{\mathcal S}}\) be given by
\[
  \widehat u(i,r):=u_i\mathbf{1}_{\{r=0\}},
  \qquad
  \widehat b(i,r):=b_r(i).
\]
Then for every \(m\geq 1\),
\[
  \varepsilon_m(\mathsf{P}_R;\nu_\varphi)
  = \widehat u^{\mathsf T}\widehat B_H^{m-1}\widehat b.
\]
Consequently the error generating function
\[
  \mathcal E_{\mathsf{P}_R}(z)
  := \sum_{m\geq 1}
      \varepsilon_m(\mathsf{P}_R;\nu_\varphi)z^m
\]
satisfies
\[
  \mathcal E_{\mathsf{P}_R}(z)
  = z\,\widehat u^{\mathsf T}
      (I-z\widehat B_H)^{-1}\widehat b
  \qquad (|z|<\rho_H^{-1}),
\]
and therefore extends to a rational function of \(z\).

If \(S_q\) denotes the cyclic permutation matrix on \(\ZZ/q\ZZ\), then
\[
  \widehat B_H=S_q\otimes B_H.
\]
Hence
\[
  \det(I-z\widehat B_H)=\det(I-z^q B_H^q),
\]
and
\[
  \spec(\widehat B_H)
  = \bigcup_{\omega^q=1}\omega\,\spec(B_H).
\]
In particular, every pole of \(\mathcal E_{\mathsf{P}_R}\) belongs to
\[
  \set{(\omega\lambda)^{-1}:
       \omega^q=1,\ \lambda\in\spec(B_H)\setminus\{0\}},
\]
and the positive real pole of smallest modulus is
\[
  z_H=\rho_H^{-1}=e^{P_{Y_H}(\varphi)-P_Y(\varphi)}.
\]
Thus finite-observation error detects the killed open-system spectrum
through a cyclotomic residue lift, while the escape rate is recovered as
the dominant real pole.
\end{theorem}

\begin{proof}
For \(n\geq 0\), a direct induction on \(n\) gives
\[
  \widehat B_H^n\bigl((i,r),(j,s)\bigr)
  = B_H^n(i,j)\mathbf{1}_{\{s=r+n\bmod q\}}.
\]
Therefore
\begin{align*}
  \widehat u^{\mathsf T}\widehat B_H^{m-1}\widehat b
  &= \sum_{i,j\in\mathcal S}
      u_i B_H^{m-1}(i,j)b_{m-1\bmod q}(j) \\
  &= u^{\mathsf T}B_H^{m-1}b_{m-1\bmod q},
\end{align*}
which is exactly
\(\varepsilon_m(\mathsf{P}_R;\nu_\varphi)\) by
Theorem~\ref{thm:first-entry-escape-rate}.

For \(|z|<\rho(\widehat B_H)^{-1}\), the geometric series converges
absolutely and yields
\[
  \mathcal E_{\mathsf{P}_R}(z)
  = \sum_{m\geq 1}
      \widehat u^{\mathsf T}\widehat B_H^{m-1}\widehat b\,z^m
  = z\,\widehat u^{\mathsf T}
      \sum_{n\geq 0}(z\widehat B_H)^n\widehat b
  = z\,\widehat u^{\mathsf T}
      (I-z\widehat B_H)^{-1}\widehat b.
\]
This is rational in \(z\).

If \(S_q\) is the cyclic permutation matrix on \(\ZZ/q\ZZ\), then the
definition of \(\widehat B_H\) gives
\[
  \widehat B_H=S_q\otimes B_H.
\]
Since the eigenvalues of \(S_q\) are exactly the \(q\)-th roots of
unity, the standard spectrum formula for Kronecker products implies
\[
  \spec(\widehat B_H)
  = \set{\omega\lambda:
         \omega^q=1,\ \lambda\in\spec(B_H)}.
\]
Multiplying over the spectrum yields
\[
  \det(I-z\widehat B_H)
  = \prod_{\lambda\in\spec(B_H)}
      \prod_{\omega^q=1}(1-z\omega\lambda)
  = \prod_{\lambda\in\spec(B_H)}(1-z^q\lambda^q)
  = \det(I-z^q B_H^q).
\]
Hence the poles of the scalar resolvent
\(\mathcal E_{\mathsf{P}_R}\) are contained in the reciprocal lifted
spectrum.

It remains to identify the dominant real pole.  Let \(v,\ell\geq 0\) be
Perron right and left eigenvectors of \(B_H\) for \(\rho_H\), normalized
by \(\ell^{\mathsf T}v=1\).  Since the transient block \(Q\) is
nilpotent, both \(v\) and \(\ell\) are strictly positive on
\(\mathcal S_\infty\).  Set
\[
  \widehat v:=\mathbf{1}_q\otimes v,
  \qquad
  \widehat\ell:=\mathbf{1}_q\otimes \ell,
\]
where \(\mathbf{1}_q\) is the all-ones vector on \(\ZZ/q\ZZ\).  Then
\(\widehat B_H\widehat v=\rho_H\widehat v\) and
\(\widehat\ell^{\mathsf T}\widehat B_H=\rho_H\widehat\ell^{\mathsf T}\).
Moreover,
\[
  \widehat u^{\mathsf T}\widehat v=u^{\mathsf T}v>0,
  \qquad
  \widehat\ell^{\mathsf T}\widehat b
  = \sum_{a\in\ZZ/q\ZZ}\ell^{\mathsf T}b_a>0,
\]
because \(u_i=\pi_i>0\) on \(\mathcal S_\infty\) and each \(b_a(i)\geq
\theta>0\).  Thus the scalar resolvent coefficient does not annihilate
the eigenprojection of \(\rho_H\), so \(z=\rho_H^{-1}\) is a genuine
simple pole of \(\mathcal E_{\mathsf{P}_R}\).  The pressure identity
from Theorem~\ref{thm:first-entry-escape-rate} gives the last formula.
\end{proof}

\begin{theorem}[H\"older Gibbs cyclotomic lift after block recoding]
\label{thm:holder-open-resonance}
Let \(Y\) be a topologically mixing one-sided subshift of finite type,
let \(\varphi\) be a H\"older potential normalized so that its transfer
operator \(\mathcal L_\varphi\) has maximal eigenvalue \(1\), and let
\(H\subset Y\) be a finite union of cylinders.  Fix \(q\geq 2\) and a
non-empty proper subset \(R\subsetneq \ZZ/q\ZZ\).  Choose
\(\ell\geq 1\) so that \(H\) is a union of \(\ell\)-cylinders, and let
\((Y^{(\ell)},\sigma^{(\ell)})\) be the \(\ell\)-block presentation of
\(Y\).  Write \(H^{(\ell)}\) for the corresponding one-step hole and
\(\varphi^{(\ell)}\) for the lifted H\"older potential.

For the first-entry residue event \(\mathsf P_R^{(\ell)}\) in the block
system, let
\(\varepsilon_n^{(\ell)}(\mathsf P_R^{(\ell)};\nu_{\varphi^{(\ell)}})\)
denote the depth-\(n\) Bayes error computed from block prefixes.  Then
the following hold.
\begin{enumerate}[label=(\roman*)]
  \item For every \(m\geq \ell\),
    \[
      \varepsilon_m(\mathsf P_R;\nu_\varphi)
      =
      \varepsilon_{m-\ell+1}^{(\ell)}
        (\mathsf P_R^{(\ell)};\nu_{\varphi^{(\ell)}}).
    \]
  \item Let \(\widehat Y^{(\ell)}\) be the natural extension of
    \(Y^{(\ell)}\), let \((Y^{(\ell)})^{-}\) be its left-history
    projection, and let \(\nu_{\varphi^{(\ell)}}^{-}\) be the past
    marginal of the equilibrium state.  The normalized potential
    \(\varphi^{(\ell)}\) determines, by the standard H\"older
    \(g\)-measure specification, a backward transfer operator
    \[
      \mathcal K_{\varphi^{(\ell)}}
      : \mathcal H_\alpha((Y^{(\ell)})^{-})
      \to \mathcal H_\alpha((Y^{(\ell)})^{-}),
    \]
    and because \(H^{(\ell)}\) is one-step, restricting to safe current
    block states gives an open backward operator
    \[
      \mathcal K_{\varphi^{(\ell)},H^{(\ell)}}.
    \]
    There exist H\"older functions
    \(g_0^-,\ldots,g_{q-1}^-\in
      \mathcal H_\alpha((Y^{(\ell)})^{-})\)
    such that, with
    \[
      \Lambda^-(f)
      := \int f\,\dd\nu_{\varphi^{(\ell)}}^{-},
    \]
    one has
    \[
      \varepsilon_n^{(\ell)}
      (\mathsf P_R^{(\ell)};\nu_{\varphi^{(\ell)}})
      =
      \Lambda^-\bigl(
        \mathcal K_{\varphi^{(\ell)},H^{(\ell)}}^{\,n-1}
        g_{n-1\bmod q}^-
      \bigr)
      \qquad (n\geq 1).
    \]
  \item Let
    \[
      \widehat{\mathcal K}F_r
      := \mathcal K_{\varphi^{(\ell)},H^{(\ell)}}(F_{r+1}),
      \qquad r\in\ZZ/q\ZZ,
    \]
    on the \(q\)-fold product Banach space, with indices read modulo
    \(q\).  Then the block-error generating function
    \[
      \mathcal E_{\mathsf P_R}^{(\ell)}(z)
      := \sum_{n\geq 1}
         \varepsilon_n^{(\ell)}
           (\mathsf P_R^{(\ell)};\nu_{\varphi^{(\ell)}})z^n
    \]
    satisfies
    \[
      \mathcal E_{\mathsf P_R}^{(\ell)}(z)
      = z\,\widehat\Lambda^-
          (I-z\widehat{\mathcal K})^{-1}\widehat g^-,
    \]
    where \(\widehat g^-:=(g_0^-,\ldots,g_{q-1}^-)\) and
    \(\widehat\Lambda^-(F):=\Lambda^-(F_0)\).  Consequently every pole
    of \(\mathcal E_{\mathsf P_R}^{(\ell)}\) belongs to the reciprocal
    \(q\)-th cyclotomic lift of
    \(\spec(\mathcal K_{\varphi^{(\ell)},H^{(\ell)}})\).
  \item The isolated spectrum of
    \(\mathcal K_{\varphi^{(\ell)},H^{(\ell)}}\) agrees, under the
    standard past--future identification coming from the natural
    extension, with the isolated spectrum of the usual forward open
    Ruelle operator \(\mathcal L_{\varphi,H}\).  Hence every pole of the
    original generating function
    \[
      \mathcal E_{\mathsf P_R}(z)
      := \sum_{m\geq 1}
         \varepsilon_m(\mathsf P_R;\nu_\varphi)z^m
    \]
    belongs to the reciprocal \(q\)-th cyclotomic lift of the open
    Ruelle spectrum.
  \item If the leading open resonance
    \(e^{-\gamma_H(\varphi)}\) of \(\mathcal L_{\varphi,H}\) is simple
    and isolated, then there exist \(\eta>0\) and positive constants
    \(c_r\) such that
    \[
      \varepsilon_m(\mathsf P_R;\nu_\varphi)
      = c_{m-\ell\bmod q}\,e^{-m\gamma_H(\varphi)}
        + O\bigl(e^{-m(\gamma_H(\varphi)+\eta)}\bigr)
      \qquad (m\to\infty).
    \]
\end{enumerate}
\end{theorem}

\begin{proof}
For part~(i), the map sending a point
\(y_0y_1y_2\cdots\in Y\) to the block sequence
\[
  (y_0\cdots y_{\ell-1})(y_1\cdots y_\ell)(y_2\cdots y_{\ell+1})\cdots
\]
is a topological conjugacy from \(Y\) onto \(Y^{(\ell)}\).  A prefix of
length \(m\) in the original coding is equivalent to a prefix of length
\(m-\ell+1\) in the block coding, and the first-entry time into \(H\)
coincides with the first-entry time into \(H^{(\ell)}\).  Hence the two
Bayes problems have identical admissible decision \(\sigma\)-algebras,
which proves the identity.

For part~(ii), pass to the block system and write
\(\psi:=\varphi^{(\ell)}\).  Since \(\psi\) is normalized and H\"older,
\(\nu_\psi\) is a H\"older \(g\)-measure in the sense of
\citet{Walters1975GMeasures}.  Therefore the natural extension
\(\widehat Y^{(\ell)}\) admits a past disintegration over the left
history space \((Y^{(\ell)})^{-}\), and the associated backward
specification operator \(\mathcal K_\psi\) preserves H\"older
regularity.  Because \(H^{(\ell)}\) is one-step, membership in the hole
depends only on the current block state, so restricting the
specification to safe current states defines the open operator
\(\mathcal K_{\psi,H^{(\ell)}}\).

For \(r\in\ZZ/q\ZZ\) and \(\xi^-\in (Y^{(\ell)})^{-}\), let
\(p_r^-(\xi^-)\) be the conditional probability, with respect to this
backward Gibbs specification, that the corresponding forward orbit
enters \(H^{(\ell)}\) for the first time in the residue class
\(R-r\).  Standard H\"older continuity of the specification kernel
implies that each \(p_r^-\) is H\"older.  Set
\[
  g_r^-:=\min\{p_r^-,1-p_r^-\}.
\]
The exact cylinder formula for Bayes error, applied on the block coding
side and then disintegrated over past histories, gives
\[
  \varepsilon_n^{(\ell)}
  (\mathsf P_R^{(\ell)};\nu_\psi)
  =
  \Lambda^-\bigl(
    \mathcal K_{\psi,H^{(\ell)}}^{\,n-1}g_{n-1\bmod q}^-
  \bigr),
\]
which is the desired representation.

For part~(iii), define \(\widehat{\mathcal K}\) on the \(q\)-fold
product exactly as in the finite-state proof.  Then
\[
  (\widehat{\mathcal K}^n\widehat g^-)_0
  = \mathcal K_{\psi,H^{(\ell)}}^{\,n}g_{n\bmod q}^-,
\]
so the geometric-series identity gives
\[
  \mathcal E_{\mathsf P_R}^{(\ell)}(z)
  = z\,\widehat\Lambda^-
      (I-z\widehat{\mathcal K})^{-1}\widehat g^-.
\]
Since \(\widehat{\mathcal K}=S_q^{-1}\otimes
\mathcal K_{\psi,H^{(\ell)}}\), the spectrum of \(\widehat{\mathcal K}\)
is the \(q\)-th cyclotomic lift of the spectrum of
\(\mathcal K_{\psi,H^{(\ell)}}\).

For part~(iv), the past and future transfer descriptions arise from the
same normalized Gibbs state on the natural extension.  Standard
past--future duality for H\"older Gibbs measures on mixing subshifts of
finite type identifies the isolated spectrum of the backward
specification operator with that of the forward Ruelle operator; see
\citet{Bowen2008,Ruelle1978,Walters1975GMeasures}.  Block recoding does
not change the open pressure or the isolated resonance set, so the same
isolated spectrum is obtained from \(\mathcal L_{\varphi,H}\).  The pole
statement for the original generating function follows from part~(i),
which shifts the index by the bounded amount \(\ell-1\).

For part~(v), write the spectral decomposition on the past Banach space
as
\[
  \mathcal K_{\psi,H^{(\ell)}}^n
  = \lambda_H^n\Pi_H^- + R_H^{-\,n},
\]
where \(\lambda_H=e^{-\gamma_H(\varphi)}\) and
\(\|R_H^{-\,n}\|=O(e^{-n(\gamma_H(\varphi)+\eta)})\).  Then
\[
  \varepsilon_n^{(\ell)}
  (\mathsf P_R^{(\ell)};\nu_\psi)
  = d_{n-1\bmod q}\,e^{-n\gamma_H(\varphi)}
    + O\bigl(e^{-n(\gamma_H(\varphi)+\eta)}\bigr)
\]
for suitable positive constants
\[
  d_r:=e^{\gamma_H(\varphi)}
       \Lambda^-(\Pi_H^- g_r^-).
\]
Use part~(i) with \(n=m-\ell+1\), and absorb the bounded shift
\(\ell-1\) into the constants to obtain
\[
  \varepsilon_m(\mathsf P_R;\nu_\varphi)
  = c_{m-\ell\bmod q}\,e^{-m\gamma_H(\varphi)}
    + O\bigl(e^{-m(\gamma_H(\varphi)+\eta)}\bigr).
\]
\end{proof}

\begin{proposition}[Doubling map and the Fibonacci escape law]
\label{prop:doubling-fibonacci-escape}
Let \(T(x)=2x\bmod 1\) on \([0,1)\), equipped with Lebesgue measure
\(\Leb\), and let \(W=[1/2,1)\).  The binary itinerary of \(x\) gives
the full one-sided two-shift with the Bernoulli \((1/2,1/2)\) measure.
Let
\[
  H=[3/4,1),
\]
so that \(H\) corresponds to the cylinder \([11]\).  Define
\[
  \tau_H(x):=\inf\set{n\geq 0:T^n x\in H}
\]
and
\[
  \mathsf{P}_{\mathrm{even}}
  := \set{x:\tau_H(x)<\infty,\ \tau_H(x)\equiv 0 \pmod 2}.
\]
If \((F_n)_{n\geq 0}\) denotes the Fibonacci sequence
\(F_0=0\), \(F_1=1\), then for every \(m\geq 1\),
\[
  \varepsilon_m(\mathsf{P}_{\mathrm{even}};\Leb)
  = \frac{F_{m+3}}{5\,2^m}.
\]
Moreover,
\[
  \Leb(\tau_H\geq m)=\frac{F_{m+3}}{2^{m+1}},
  \qquad
  \varepsilon_m(\mathsf{P}_{\mathrm{even}};\Leb)
  = \frac25\,\Leb(\tau_H\geq m),
\]
and therefore
\[
  \lim_{m\to\infty}
  -\frac1m\log\varepsilon_m(\mathsf{P}_{\mathrm{even}};\Leb)
  = \log\frac{2}{\varphi}
  = \log 2-\log\varphi.
\]
Equivalently, the survivor entropy of the golden-mean subshift is
recovered from finite-observation error by
\[
  h_{\mathrm{top}}(Y_H)
  = \log\varphi
  = \log 2
    + \lim_{m\to\infty}\frac1m
      \log \varepsilon_m(\mathsf{P}_{\mathrm{even}};\Leb).
\]
\end{proposition}

\begin{proof}
In binary coding, \(\tau_H\) is the first index at which the block
\(11\) appears.  A length-\(m\) prefix contributes to
\(\varepsilon_m(\mathsf{P}_{\mathrm{even}};\Leb)\) if and only if it
contains no occurrence of \(11\) at positions \(0,\ldots,m-2\).  Let
\begin{align*}
  G_m^0
  &:= \set{a\in\{0,1\}^m:
           a \text{ avoids }11,\ a_{m-1}=0}, \\
  G_m^1
  &:= \set{a\in\{0,1\}^m:
           a \text{ avoids }11,\ a_{m-1}=1}.
\end{align*}
The usual golden-mean recursion gives
\[
  \card{G_m^0}=F_{m+1},
  \qquad
  \card{G_m^1}=F_m.
\]
Every length-\(m\) cylinder has Lebesgue measure \(2^{-m}\).

Let
\begin{align*}
  a
  &:= \Leb(\mathsf{P}_{\mathrm{even}}
           \mid \text{safe prefix of even length ending in }0), \\
  b
  &:= \Leb(\mathsf{P}_{\mathrm{even}}
           \mid \text{safe prefix of even length ending in }1), \\
  c
  &:= \Leb(\mathsf{P}_{\mathrm{even}}
           \mid \text{safe prefix of odd length ending in }0), \\
  d
  &:= \Leb(\mathsf{P}_{\mathrm{even}}
           \mid \text{safe prefix of odd length ending in }1).
\end{align*}
Since the next binary digit is \(0\) or \(1\) with equal probability,
\[
  a=\frac12 c+\frac12 d,
  \qquad
  b=\frac12 c,
\]
\[
  c=\frac12 a+\frac12 b,
  \qquad
  d=\frac12+\frac12 a.
\]
Solving gives
\[
  a=\frac35,\qquad b=\frac15,\qquad c=\frac25,\qquad d=\frac45.
\]
Hence the Bayes error on a safe prefix depends only on its terminal
symbol:
\[
  \beta(a)=\frac25
  \quad\text{on }G_m^0,
  \qquad
  \beta(a)=\frac15
  \quad\text{on }G_m^1.
\]
Therefore
\[
  \varepsilon_m(\mathsf{P}_{\mathrm{even}};\Leb)
  = 2^{-m}\left(
      \frac25\card{G_m^0}
      + \frac15\card{G_m^1}
    \right)
  = 2^{-m}\left(
      \frac25F_{m+1}
      + \frac15F_m
    \right).
\]
Using \(F_{m+3}=2F_{m+1}+F_m\), we obtain
\[
  \varepsilon_m(\mathsf{P}_{\mathrm{even}};\Leb)
  = \frac{F_{m+3}}{5\,2^m}.
\]

The condition \(\tau_H\geq m\) is equivalent to saying that the first
\(m+1\) binary digits avoid \(11\), so
\[
  \Leb(\tau_H\geq m)=\frac{F_{m+3}}{2^{m+1}}.
\]
This proves the exact factor \(\varepsilon_m=\frac25\Leb(\tau_H\geq m)\).
Finally, \(F_{m+3}\sim \varphi^{m+3}/\sqrt5\), so
\[
  \lim_{m\to\infty}
  -\frac1m\log\varepsilon_m(\mathsf{P}_{\mathrm{even}};\Leb)
  = \log\frac{2}{\varphi}.
\]
The survivor subshift is the golden-mean shift, whose topological
entropy is \(\log\varphi\), and the last identity follows.
\end{proof}

\begin{proof}[Proof of Theorem~\ref{thm:main-resonance}]
The exact Markov resolvent statement is
Theorem~\ref{thm:error-resolvent-cyclotomic}.  The H\"older Gibbs
extension, under the transfer-operator representation and spectral-gap
hypotheses, is Theorem~\ref{thm:holder-open-resonance}.
\end{proof}
